% !TEX root = ../main.tex

% \begin{abstract}[zh]

%   粒子物理致力于研究组成宇宙的基本粒子及其相互作用，其中的标准模型是目前最为成功的理论框架，将所有已知的基本粒子和三种基本相互作用统一描述，并成功地由大量实验验证。
%   然而，实验上的许多现象如暗物质、暗能量、物质-反物质不对称性等尚无法由标准模型解释，模型本身也存在如层级问题、强CP问题等理论上的不足。
%   许多超出标准模型的理论试图解决这些问题，其中包括双希格斯双重态模型和有效场论等。同时，这些理论预言了新粒子与标准模型粒子之间的相互作用，为实验探索新物理提供了重要途径。
%   顶夸克作为标准模型中质量最大的基本粒子，其与希格斯场的强耦合使其成为新物理研究的关键探针。
%   四顶夸克产生过程是标准模型中极为罕见的反应之一，并于2023年首次被ATLAS实验观测，其测量结果在不确定度范围内与标准模型预言一致，但中心值的偏高仍为新物理效应留下了潜在空间。

%   本论文基于大型强子对撞机ATLAS实验在$\sqrt{s}=13~\mathrm{TeV}$下收集的139~fb$^{-1}$ Run~2数据，在单轻子及双反电荷轻子末态中开展四顶夸克事例的研究，
%   重点搜索双希格斯双重态模型中的重标量粒子（$H/A$）以及色八重态标量粒子（$S_8$）所产生的四顶夸克信号。
%   针对主导背景$t\bar{t}+\text{jets}$，采用基于数据驱动的味成分重缩放方法以及基于神经网络的多维动力学重加权进行建模修正。
%   同时，引入质量参数化的图神经网络以提升信号与背景的区分能力，并与多轻子末态分析结果进行联合拟合。
%   在95\%置信水平下，对于质量为400~GeV和1000~GeV的$H/A$，分别得到$t\bar{t}H/A \to t\bar{t}t\bar{t}$过程的截面上限为14~fb和5.0~fb。
%   在$H$与$A$共同贡献的假设下，对应排除$\tan\beta < 1.7$和$0.7$的参数空间。同时，在$S_8S_8 \to t\bar{t}t\bar{t}$的搜索中排除了质量低于1.3~TeV的$S_8$。
%   基于高亮度大型强子对撞机（HL-LHC）的展望研究表明，在$3~\mathrm{ab}^{-1}$、$\sqrt{s}=14~\mathrm{TeV}$条件下，对于质量超过1~TeV的新粒子，其截面灵敏度可提升至1~fb以下。

%   为进一步提升四顶夸克过程的物理测量能力，本文构建了适用于Run~3数据的微分截面测量框架。
%   该框架基于图神经网络实现四顶夸克衰变链的重建，并结合解析方法与生成模型对中微子动量进行推断，从而重建部分子层级观测量，如四顶夸克不变质量$m_{4t}$。
%   在此基础上，采用基于剖面似然函数的展开方法提取微分截面，并系统研究了正则化策略及方法稳健性。该框架为未来四顶夸克微分测量提供了系统的方法学基础。

%   此外，本文还开展了与ATLAS阻性板室（RPC）探测器相关的研究工作。包括基于探测器控制系统数据开发自动化性能监控工具，以及参与Phase-II升级相关的读出板设计与质量控制测试研究。
%   同时，通过有限元方法对气体流动与电场分布进行模拟，评估结构参数对探测器性能的影响，为未来优化设计提供参考。

% \end{abstract}

\begin{abstract}[en]

  Particle physics aims to understand the fundamental constituents of matter and their interactions. 
  The Standard Model (SM) provides a remarkably successful theoretical framework, describing all known elementary particles and three of the four fundamental forces, and has been extensively validated by experimental measurements. 
  However, several observations, such as dark matter, dark energy, and the matter--antimatter asymmetry of the Universe, remain unexplained within the SM. 
  In addition, the SM faces theoretical challenges, including the hierarchy problem and the strong CP problem. 
  Various extensions of the SM, such as Two-Higgs-Doublet Models (2HDM) and Effective Field Theory (EFT) approaches, have been proposed to address these issues, often predicting new particles with enhanced couplings to the top quark.

  As the heaviest known elementary particle, the top quark plays a central role in many beyond-the-Standard-Model scenarios due to its strong coupling to the Higgs sector. 
  The production of four top quarks is a particularly rare process in the SM and was observed for the first time by the ATLAS experiment in 2023. 
  While the measured cross section is consistent with the SM prediction within uncertainties, the central value remains slightly elevated, leaving room for potential contributions from new physics.

  This thesis presents a search for heavy resonance in $t\bar{t}t\bar{t}$ final states using events with one lepton or two opposite-sign leptons and high jet multiplicities. 
  The analysis is based on the full Run~2 dataset collected by the ATLAS experiment at $\sqrt{s}=13~\mathrm{TeV}$, corresponding to an integrated luminosity of 139~fb$^{-1}$. 
  The search targets heavy scalar and pseudoscalar particles ($H/A$) in the context of 2HDM, produced in association with a top-quark pair and decaying into $t\bar{t}$, as well as pair-produced colour-octet scalars ($S_8$). 
  The dominant $t\bar{t}+\mathrm{jets}$ background is controlled using a data-driven flavour-dependent normalisation and a neural-network-based multidimensional kinematic reweighting. 
  An $H/A$-mass parameterised Graph Neural Network is employed to enhance the discrimination between signal and background, and the results are combined with two same-sign leptons or multilepton analyses.
  No significant deviation from the SM prediction is observed. 
  Upper limits at 95\% confidence level are set on 
  $\sigma(pp \to t\bar{t}H/A)\times \mathcal{B}(H/A \to t\bar{t})$, 
  with values of 14~fb (5.0~fb) at 400~GeV (1000~GeV), improving upon the same-sign and multilepton channels by up to 19\%.
  In the scenario where both $H$ and $A$ contribute, values of $\tan\beta$ below 1.7 and 0.7 are excluded for masses of 400~GeV and 1000~GeV, respectively. 
  For $S_8$ pair production, masses below 1.3~TeV are excluded. 
  Projections to the High-Luminosity LHC (HL-LHC) indicate that, with $3~\mathrm{ab}^{-1}$ of data at $\sqrt{s}=14~\mathrm{TeV}$, cross-section sensitivities below 1~fb can be achieved for heavy resonance masses above 1~TeV.

  To further exploit the physics potential of $t\bar{t}t\bar{t}$ production, a reconstruction and unfolding framework is developed for differential measurements in Run~3 data. 
  A dedicated reconstruction strategy based on the blended graph-hypergraph neural network framework is introduced to assign reconstructed objects to the four top-quark decay chains. 
  The neutrino kinematics are inferred using both analytic techniques and machine-learning-based approaches, enabling the reconstruction of parton-level observables such as the $t\bar{t}t\bar{t}$ invariant mass $m_{4t}$. 
  Differential cross sections are extracted using a profile likelihood unfolding approach, with regularisation applied to ensure stability in the presence of limited statistics. 
  Extensive validation studies demonstrate the robustness and reliability of the method.

  In addition, contributions to the ATLAS muon spectrometer are presented through work on Resistive Plate Chambers (RPCs). 
  These include the development of automated monitoring tools based on Detector Control System data, as well as studies in the context of the Phase-II upgrade, focusing on readout board production and quality control, and mechanical tests of RPC gas gaps. 
  Separately, dedicated studies are carried out using finite-element simulations of gas flow, electric field distributions, and electrode deformation to assess the impact of structural configurations and guide future detector optimisation.
 
  % Overall, this thesis demonstrates the strong potential of the $t\bar{t}t\bar{t}$ topology as a probe of both the Standard Model and new physics. 
  % The combination of advanced analysis techniques, improved reconstruction methodologies, and detector developments provides a comprehensive framework for future studies at the LHC and the HL-LHC.

\end{abstract}
